\section{Un orden común para todos los menores reales}
Las claves de las fuentes se desarrollan en el catálogo literal. U01--05 y F02--03
usan dos enteros distintos: $j$ fija el orden y $m$ fija el punto de evaluación.
Sea $D$ bistocástica, $jD_{rc}\le M_0$ y $\|D-J_j\|_2\le r<1$ (norma de operador).
Se puede formar primero el exterior común de U02 borrando $v=o(j)$ sitios.
En el gas global basta $v=0$. Sean $m=j-h$, $0\le h\le a_*j$ y
$Q_*=\lfloor\varepsilon j\rfloor$, con $a_*,\varepsilon>0$ de clase. Distinguimos
este orden del conteo $Q(B)$. U02 da en cada menor su propio escalado $D'$,
en el mismo soporte, $mD'_{rc}\le C$ y $\|B\|_2\le r'<1$, $B=D'-J_m$.
Los argumentos que siguen también valen con $h=0$.

\paragraph{Polinomio y disco físico.}
Pongamos $R_B(z)=\sum_S z^{|S|}\prod_{e\in S}B_e$, sobre matchings, incluido el vacío.
En la descomposición singular real $B=U\Sigma V^t$ tomamos sólo los valores positivos.
Para gaussianas circulares estándar independientes $Z_l$, de segundo momento uno, sea
\[
 u_i=\sum_l U_{il}\sqrt{\sigma_l} Z_l,\qquad
 v_j=\sum_l V_{jl}\sqrt{\sigma_l}\overline Z_l,\qquad W=\sum_l\sigma_l|Z_l|^2.
\]
Las dos sumas de coordenadas son cero. En cada lado la suma de módulos cuadrados es $W$.
Wick da $R_B(z)=\mathbb E\prod_i(1+\sqrt z\,u_i)\prod_l(1+\sqrt z\,v_l)$:
cada emparejamiento cruza los lados y lleva un factor $zB_{il}$.
La media geométrica de los cuadrados de los módulos de cada lado da
\[
 |R_B(z)|\le\mathbb E(1+|z|W/m)^m .
\]
En $|z|r'<1$ con margen, $\log\mathbb E e^{|z|W}\le C\sum\sigma_l\le C\sqrt m$
por difusión y Cauchy--Schwarz. Sea $d(z)=\det(I-z^2B^tB)^{-1/2}$, rama desde uno.
$d^{\pm1}$ están acotados en el mismo disco reducido por $\sum\sigma_l^2\le C$.
Fijamos $r_1>1$ y una ventana $|y-1|<d_0$ con $r_1(1+d_0)r'<1$ con holgura.
Se obtiene el mayorante $e^{C\sqrt m+C}$ en un disco de $w$ algo mayor que $r_1$
para $F_y(w)=R_B(yw)/d(yw)$.

\paragraph{Identidad formal y truncación.}
Para el vector gaussiano bilineal formal de covarianza
$O=\left(\begin{smallmatrix}0&ywB\\ywB^t&0\end{smallmatrix}\right)$,
extraer $e^{-\sum_i X_i^2/2}$ del producto deja los factores
$(1+X_i)e^{-X_i+X_i^2/2}$, pues $\sum X_i=0$.
Completando el cuadrado formalmente se obtiene $d(yw)$ y la covarianza $O(I+O)^{-1}$.
Esto también se verifica coeficiente a coeficiente con una fuente lineal: la razón de las
dos funciones generatrices tras el giro cuadrático es
$\exp(t^tO(I+O)^{-1}t/2)$. Se aplica cerca de cero y después sólo como identidad formal.
Cada vértice activo tiene grado al menos tres y cada arista empieza en $w$.
Con $E$ aristas y $s$ vértices, $l=E-s$ implica $s\le2l$, $E\le3l$.
U03, líneas 9--43, da la fórmula explícita $V_l$ por inversión sobre índices coincidentes.
Su factor combinatorio se acota por $(2E-1)!!/s!\le C^E E^{E-s}$ y
$\prod_{i<s}(1+i|u|)\le e^{Cs}$ para $l\le Q_*$, $|u|\le c_0/Q_*$.
El núcleo normalizado de cada arista es acotado por U02. Resulta
$|V_l|\le(Cl|u|)^l$. En las medidas físicas y $u=1/m$ la inversión es exacta.
Hasta grado $Q_*$ en $w$ no cabe ningún diagrama omitido, ya que $E\le Q_*$ allí.
Cauchy al evaluar en $w=1$ acota la diferencia entre $R_B(y)$ y
$S_{Q_*}(y,1/m)=d(y)(1+\sum_{l\le Q_*}V_l)$ por
$C(e^{C\sqrt m}+1)r_1^{-Q_*}=O(e^{-cj})$. Aquí y abajo $j$ es suficientemente grande.

\paragraph{Cola ponderada e integración.}
Para $T\sim\Gamma(m+1,1)$ sea $X=m/T$. Expandir $\operatorname{per}(J_m+B)$
y usar $\mathbb E X^i=m^i(m-i)!/m!$ da
$m^m\operatorname{per}D'/m!=\mathbb E R_B(X)$ exactamente.
En $X\notin I=[1-d,1+d]$ el mayorante anterior cancela el $T^m$ de la densidad.
El cambio $t'=t+W$ deja una cola gamma de región
$\{t:m/t\notin I\}+W$, con peso $e^W$.
Si $W\le\kappa_\Gamma m$, para $\kappa_\Gamma$ pequeño las dos partes quedan a distancia
positiva proporcional a $m$ de la media. Chernoff da $Ce^{-cm}$; el peso satisface
$\mathbb E e^W\le e^{C\sqrt m}$.
En $W>\kappa_\Gamma m$ usamos $e^{-\gamma\kappa_\Gamma m}\mathbb E e^{(1+\gamma)W}$
con $(1+\gamma)r'<1$. La cola ponderada es $O(e^{-cm+C\sqrt m})$.

Sea $\mu_i(u)$ la recurrencia de U04, líneas 19--27. A $u=1/m$ son los momentos
centrados de $X$ en uno. Para $N|u|$ pequeño, $t=N|u|/(1-N|u|)$,
la inducción en la recurrencia da $|\mu_i|\le(2\sqrt t)^i$, $i\le N$.
Fijamos $d<r_0<d_0$. Para $2Q_*+2$ el momento par paga el resto de Taylor en $I$
por $(CQ_*/m)^{Q_*+1}$. En $I^c$ cada monomio $|X-1|^i$, $i\le2Q_*+1$,
se acota por $d^{i-2Q_*-2}|X-1|^{2Q_*+2}$ y se suma con $r_0^{-i}$.
En particular es lícito sustituir la integral por
\[
 H_{Q_*}(u)=\sum_{i=0}^{2Q_*+1}\mu_i(u)\,\partial_y^i S_{Q_*}(1,u)/i!.
\]
Para $|u|\le c_0/Q_*$, $\sum_l(Cl|u|)^l=O(|u|)$, y tomar $N=i$
en la recurrencia da $\sum_{i\ge2}^{2Q_*+1}|\mu_i|r_0^{-i}=O(|u|)$.
De aquí $H_{Q_*}=d(1)+O(|u|)$, separado de cero, con rama logarítmica acotada.
Reducimos $\varepsilon$ para todos los $m\ge(1-a_*)j$. La comparación real es
\[
 \left|\log(m^m\operatorname{per}D'/m!)-\log H_{Q_*}(1/m;D')\right|\le Ce^{-cj}.
\]
La parte formal, el disco físico y la cola ponderada desempeñan tres funciones distintas.
En datos complejos se prolonga el aproximante definido por U02--04.
